%%%%%%%%%%%%%%%%%%%%%%%%%%%%%%%%%%%%%%%%%%%%%%%%%%%%%%%%%%%%%%%%%%%%%%%%%%%%%%%%
%%  EINLEITUNG  (wird ZULETZT geschrieben)   Umfang: 2–4 Seiten
%%  Gewicht 15 % der Note. Grundlage: dein Motivationsformular (Anhang).
%%%%%%%%%%%%%%%%%%%%%%%%%%%%%%%%%%%%%%%%%%%%%%%%%%%%%%%%%%%%%%%%%%%%%%%%%%%%%%%%
\chapter{Einleitung}
\label{chap:einleitung}

\begin{hinweise}
In diesem Abschnitt erörtern Sie die folgenden Themen (und nicht nur in einem Satz!):
\begin{itemize}\itemsep0pt
\item Motivation für die Wahl Ihres Projektthemas.
\item Motivation für die Wahl Ihres Schwerpunktthemas innerhalb Ihres Projektthemas.
\item Wie unterscheidet sich Ihr Projektthema von anderen Themen im aktuellen Semester?
\item Wie unterscheidet sich Ihr Projektthema von ähnlichen Themen in früheren Jahren?
\item Wie unterscheidet sich Ihr Schwerpunktthema von dem des anderen Gruppenmitglieds? (wenn Sie eine Zweiergruppe bilden)
\item Warum wollen Sie das von Ihnen gewählte Schwerpunktthema bearbeiten?
\item Was ist der Kontext Ihres Schwerpunktthemas in Ihrem Projekt?
\item Welche Anregungen hatten Sie?
\item Was erwarten Sie von der Bearbeitung Ihres Schwerpunktthemas?
\item Wie steht es mit Ihren Vorkenntnissen in Bezug zum Projektthema?
\item Wie steht es mit Ihren Vorkenntnissen in Bezug zum Schwerpunktthema?
\item Gibt es weitere Punkte, die Sie zu Ihren Vorkenntnissen anmerken wollen?
\item Weitere für Sie wichtige Punkte zu diesem Portfolio.
\item Vorschau über das Portfolio bestehend aus einer Auflistung der einzelnen Bestandteile sowie einer kurze Einführung in deren Inhalte.
\end{itemize}
Sie können dazu das Motivationsformular (als Anlage beilegen) als Grundlage verwenden.
\textbf{Die Einleitung ist mindestens zwei Seiten und maximal vier Seiten lang. Die Längenangaben sind ohne Titel, Bilder, Tabellen und Tabellen- oder Bildlegenden zu verstehen.}

\smallskip\emph{Zusätzlich aus dem Leitfaden \enquote{Portfolios im Projektpraktikum} und den Bewertungskriterien:}
Welche Gründe aus Ihrer Lernbiografie haben Sie zu dem Thema geführt? Zusammenhang Schwerpunktthema
\(\leftrightarrow\) Projektthema in mehr als einem Satz. Schwerpunkt von dem der Partnerin/des Partners abgrenzen:
\emph{Was außer dem Titel ist der Unterschied?} Das Schwerpunktthema muss klar definiert sein (fehlt es, gibt es Abzug).
Tipp: Unterabschnitte (\texttt{\textbackslash section*}) helfen, die geforderten Aussagen identifizierbar zu machen.
\end{hinweise}

\begin{umfang}{Einleitung}{2}{4}

\section*{Motivation und Themenwahl}
\platzhalter{Motivation für das Projektthema; Anregungen; Lernbiografie; Vorkenntnisse}

\section*{Schwerpunktthema und Abgrenzung}
\platzhalter{Definition des Schwerpunkts \Schwerpunkt; Kontext im Projekt; Abgrenzung zu
\Partner{} (\PartnerSchwerpunkt); Abgrenzung zu anderen/früheren Themen; Erwartungen}

\section*{Vorschau auf das Portfolio}
\platzhalter{Auflistung aller Bestandteile mit kurzer Einführung in die Inhalte, z.B. mit Verweisen
wie Abschnitt~\ref{sec:arbeitsplan}}

\end{umfang}
